\chapter{參考系、光與時空幾何}

\section{課程再次從相對論開始}

本學季回到經典物理，但不是回到空間與時間彼此獨立且普適的牛頓圖像。本學季直接探討狹義相對論與經典場論，尤其是電磁學，不會用到量子力學。先前的力學課程介紹了作用量原理，以及哈密頓和拉格朗日描述；量子課程則介紹了量子力學的基本原理，不過連諧振子也只好留待日後再談。先前也幾乎沒有談到光。現在要補上這個缺漏，因為光本質上是相對論性的。

班上大多數人以前都學過狹義相對論。儘管如此，舉手表決後，大家傾向於從頭重新建構一次。這是很好的選擇。介紹洛侖茲變換時，不該把它當成需要背誦的公式；它應該從測量空間與時間的操作程序中推導出來，因為這些測量不可能同時符合伽利略運動學與觀測到的光傳播定律。

\subsection{座標是物理測量工具}

先從笛卡兒座標系開始：一個原點、三條有方向的軸，以及座標\(x,y,z\)。具體來說，可以用剛性的米尺組成點陣，遍布整個空間。要標示事件，而不只是地點，點陣上的每個位置還都要有一個時鐘。這樣，事件就能標記為
\begin{equation}
(x,y,z,t).
\end{equation}
這些時鐘必須依照明確說明的程序同步；稍後會看到，這項要求正是整個建構中最微妙的部分。

\begin{definition}
參考系由空間座標點陣和一組分布於各處、彼此同步的時鐘構成。慣性參考系是不旋轉、不加速的參考系；在其中，每個不受力的粒子都以定速度運動。
\end{definition}

兩個彼此做相對等速運動的參考系，會為同一事件指定不同座標。固定在運動參考系中的點，其座標\(x'\)保持不變，即使它的不帶撇座標持續改變。在牛頓物理中，沿\(x\)方向運動時，空間座標的關係是
\begin{equation}
x'=x-vt.
\label{eq:sr1-galilean-space}
\end{equation}
較不顯眼的牛頓假設是：所有校準妥當的時鐘都能普遍保持同步：
\begin{equation}
t'=t.
\label{eq:sr1-absolute-time}
\end{equation}

\subsection{相對性原理}

相對論並非始於愛因斯坦。伽利略和牛頓早已採用這項原理：物理定律在每個慣性參考系中的形式都相同。牛頓定律\(F=ma\)、平方反比引力定律和庫侖力，不會只因實驗室做等速運動，就改變其形式。

常見的檢驗方式，是想像一節密閉的火車車廂。如果列車平穩行駛，窗戶又看不到外界，人在車廂內拋接球的情形就和在靜止的房間裡一樣。球照常上升、落下，再回到手中。只在車廂內進行的實驗，無法測出車廂的等速運動。加速度或轉動可以在車廂內部測出來；等速平移則不行。

愛因斯坦在這項古老原理之外，再加上一條物理定律：光在真空中的傳播速率為
\begin{equation}
c \simeq 1.86\times 10^5\ {\rm miles/s}
   \simeq 3.00\times 10^8\ {\rm m/s}.
\end{equation}
如果這是自然定律，而所有慣性參考系中的定律形式都相同，那麼每位慣性觀測者測得的\(c\)都必須相同。

\section{伽利略變換為何行不通}

只要考慮一個空間維度，就足以看出衝突。畫一張時空圖，以\(x\)為水平座標、\(t\)為垂直座標。從原點發出、向右傳播的光脈衝滿足
\begin{equation}
x=ct.
\label{eq:sr1-light-right}
\end{equation}
有一個參考系以速率\(v\)向右運動，它的原點位於
\begin{equation}
x=vt,
\qquad
x'=0.
\label{eq:sr1-moving-origin}
\end{equation}

\begin{figure}[htbp]
\centering
\includegraphics[width=0.90\textwidth]{lecture_01_figure_01.png}
\caption{向右傳播的光線與運動中的原點\(x=vt\)；帶撇座標的觀測者將該原點標為\(x'=0\)。}
\lectureframe{00:15:24}
\end{figure}

將伽利略規則\eqref{eq:sr1-galilean-space}--\eqref{eq:sr1-absolute-time}套用到光線上：
\begin{align}
x'&=x-vt=(c-v)t,\\
t'&=t.
\end{align}
因此，帶撇座標的觀測者會推得
\begin{equation}
x'=(c-v)t',
\end{equation}
也就是說，向右傳播的光速會是\(c-v\)。對於向左傳播的光線\(x=-ct\)，則有
\begin{equation}
x'=-(c+v)t',
\end{equation}
其速率大小為\(c+v\)。

\begin{figure}[htbp]
\centering
\includegraphics[width=0.90\textwidth]{lecture_01_figure_02.png}
\caption{伽利略代換預測向右傳播的光速為\(c-v\)，這與光速不因參考系而異相矛盾。}
\lectureframe{00:17:23}
\end{figure}

普通的速度相加法則，對緩慢的機械運動而言是很好的近似；但測量光速時，並未觀察到預測的\(c-v\)與\(c+v\)。需要重新檢討的是變換定律，而不是光速公設。最可疑的，就是看似無害的陳述\(t'=t\)：它主張存在絕對的同時性。

\section{如何同步相距遙遠的時鐘}

要判定相距不同位置的兩個時鐘是否顯示相同時間，必須進行實驗。在兩個時鐘的中點安排第三位觀測者，並用米尺確認兩側距離相等。請兩端的時鐘在各自顯示正午時發出閃光。如果兩道閃光同時抵達中點，就稱這兩個時鐘已同步。

這不是套用在物理之上的哲學定義，而是運用時鐘、距離和光訊號所做的可重複操作。

現在讓一位觀測者經過這個中點。原本同時抵達靜止中點的兩端閃光，不會同時抵達這位運動中的觀測者：光在傳播途中，觀測者會朝其中一道閃光移動，同時遠離另一道閃光。因此，靜止時鐘判定為同時的事件，按照運動參考系中同樣有效的同步程序來判定，就不是同時發生。

在不帶撇座標的參考系中，同時發生的一系列事件構成一條直線：
\begin{equation}
t=\text{constant}.
\end{equation}
在\(x\)-\(t\)圖中，相應的帶撇座標直線會傾斜。在推導變換公式之前，我們可以先利用光訊號建構這個傾斜。

\section{以光訊號描繪運動參考系}

令\(c=1\)。若時間以秒為單位，可以用光秒測量距離；或者選定一個長度單位，再用光穿越該距離所需的時間來測量時間。無論採用哪種約定，光線方程都是
\begin{equation}
x=\pm t,
\label{eq:sr1-light-units}
\end{equation}
因此，在兩軸比例相同的圖上，光線都是\(45^\circ\)直線。

先畫出光線，再畫出一群靜止於運動參考系中的觀測者：
\begin{equation}
x=vt,\qquad x=vt+1,\qquad x=vt+2.
\label{eq:sr1-moving-family}
\end{equation}
由於這些觀測者在自己的慣性座標點陣中保持固定間距，他們的世界線彼此平行。

\begin{figure}[htbp]
\centering
\includegraphics[width=0.90\textwidth]{lecture_01_figure_03.png}
\caption{位於\(45^\circ\)角度的光線，以及運動座標點陣中彼此平行的世界線\(x=vt\)、\(x=vt+1\)和\(x=vt+2\)。}
\lectureframe{00:29:55}
\end{figure}

\subsection{Fred、Mary 和 Seymour}

令 Fred 位於\(x=vt\)，Mary 位於\(x=vt+1\)，Seymour 位於\(x=vt+2\)。在原點\(O\)，Fred 向 Mary 發出一道向右傳播的光訊號。Seymour 也發出一道向左傳播的光訊號，讓兩道閃光在同一事件\(a\)抵達 Mary。按照運動參考系的中點同步規則，Seymour 發出訊號的事件\(b\)，與 Fred 在\(O\)發出訊號的事件同時。因此，通過\(O\)和\(b\)的直線就是\(t'=0\)。

事件\(a\)同時位於
\begin{equation}
x=t
\end{equation}
以及
\begin{equation}
x=vt+1.
\end{equation}
上。因此，
\begin{equation}
t_a=x_a=\frac{1}{1-v}.
\label{eq:sr1-event-a}
\end{equation}

通過\(a\)、向左傳播的光線滿足
\begin{equation}
x+t=x_a+t_a=\frac{2}{1-v}.
\label{eq:sr1-ab-light}
\end{equation}
事件\(b\)也位於 Seymour 的世界線上，
\begin{equation}
x-vt=2.
\label{eq:sr1-seymour-line}
\end{equation}
聯立\eqref{eq:sr1-ab-light}與\eqref{eq:sr1-seymour-line}可得
\begin{align}
t_b&=\frac{2v}{1-v^2},\\
x_b&=\frac{2}{1-v^2}.
\label{eq:sr1-event-b}
\end{align}
因此，從\(O\)到\(b\)的直線斜率為
\begin{equation}
\frac{t_b}{x_b}=v,
\end{equation}
所以運動參考系中的同時線是
\begin{equation}
t=vx,
\qquad
t'=0.
\label{eq:sr1-moving-simultaneity}
\end{equation}

\begin{figure}[htbp]
\centering
\includegraphics[width=0.93\textwidth]{lecture_01_figure_04.png}
\caption{Fred--Mary--Seymour 同步構造。光訊號同時抵達，使運動參考系將\(O\)和\(b\)判定為同時事件，因而得到\(t=vx\)。}
\lectureframe{00:44:38}
\end{figure}

現在，帶撇座標參考系的兩條座標軸可分別寫成：
\begin{align}
x'=0 &\quad\Longleftrightarrow\quad x=vt,\\
t'=0 &\quad\Longleftrightarrow\quad t=vx.
\end{align}
它們以向右傳播的光線\(x=t\)為對稱軸，彼此鏡射。幾何關係已取代了普遍適用的同時性。

\begin{classroomqa}
\audiencequestion{當\(v=1\)或\(v=-1\)時，這個建構會如何變化？}\lecturerresponse{在\eqref{eq:sr1-event-a}--\eqref{eq:sr1-event-b}中，分母會變成零，相關交點則延伸至無窮遠。當\(\lvert v\rvert=1\)時，物質觀測者會沿著光線運動，無法完成一般的中點訊號交換。這種奇異行為是第一個幾何上的警訊，顯示\(c\)是極限速率。\lecturetimestamp{00:40:43}}\end{classroomqa}

\section{推導洛侖茲變換}

在\(x=vt\)上，\(x'=0\)必須成立，因此\(x'\)必須含有因子\(x-vt\)。同樣地，在\(t=vx\)上，\(t'=0\)必須成立，因此也必須含有因子\(t-vx\)。時空具有均勻性，所以變換是線性的，可寫成
\begin{align}
x'&=(x-vt)f(v),\\
t'&=(t-vx)g(v).
\label{eq:sr1-linear-ansatz}
\end{align}

\begin{figure}[htbp]
\centering
\includegraphics[width=0.90\textwidth]{lecture_01_figure_05.png}
\caption{這個線性變換根據兩個為零的條件建構而成，並包含依速度而變的比例因子\(f(v)\)和\(g(v)\)。}
\lectureframe{00:49:35}
\end{figure}

沿著向右傳播的光線\(x=t\)，光速不變要求\(x'=t'\)。因此，方程~\eqref{eq:sr1-linear-ansatz}化為
\begin{align}
x'&=(1-v)x f(v),\\
t'&=(1-v)x g(v),
\end{align}
於是
\begin{equation}
f(v)=g(v).
\label{eq:sr1-f-equals-g}
\end{equation}
兩個方程因而共用同一個比例因子：
\begin{align}
x'&=(x-vt)f(v),\\
t'&=(t-vx)f(v).
\label{eq:sr1-direct-unscaled}
\end{align}
這基本上就是愛因斯坦第一篇相對論論文採用的簡潔論證：操作性定義的同時性與不變光速共同決定了變換。

\subsection{互易性決定比例因子}

沒有哪個慣性參考系是特別的、\emph{真正運動著}的參考系。如果帶撇參考系相對於不帶撇參考系以\(v\)運動，那麼不帶撇參考系相對於帶撇參考系就以\(-v\)運動。因此，逆變換的形式必須相同，只要將\(v\)的符號反轉：
\begin{align}
x&=(x'+vt')f(v),\\
t&=(t'+vx')f(v).
\label{eq:sr1-inverse-unscaled}
\end{align}

\begin{figure}[htbp]
\centering
\includegraphics[width=0.90\textwidth]{lecture_01_figure_06.png}
\caption{將相對速度反向並保留相同的比例因子，就能依互易性得到逆變換。}
\lectureframe{00:55:12}
\end{figure}

將\eqref{eq:sr1-direct-unscaled}代入\eqref{eq:sr1-inverse-unscaled}的第一個方程：
\begin{align}
x
&=\left[(x-vt)f(v)+v(t-vx)f(v)\right]f(v)\\
&=x(1-v^2)f(v)^2.
\end{align}
由於這兩次變換的合成必須是恆等變換，
\begin{equation}
(1-v^2)f(v)^2=1,
\end{equation}
因此
\begin{equation}
f(v)=\pm\frac{1}{\sqrt{1-v^2}}.
\label{eq:sr1-factor-roots}
\end{equation}

\begin{figure}[htbp]
\centering
\includegraphics[width=0.90\textwidth]{lecture_01_figure_07.png}
\caption{將正向變換與逆變換合成，就能確定共同因子\(f(v)\)的大小。}
\lectureframe{00:59:31}
\end{figure}

\begin{classroomqa}
\audiencequestion{為什麼在\eqref{eq:sr1-factor-roots}中要捨去負平方根？}\lecturerresponse{相對速度為零時，變換必須連續趨近恆等變換，也就是\(x'=x\)和\(t'=t\)。負號分支則給出\(x'=-x\)和\(t'=-t\)，使兩個座標軸的方向都反轉。選擇負號分支會導致突如其來的鏡射，而不是連續變化的參考系。\lecturetimestamp{00:59:27}}\end{classroomqa}

取正號分支，在\(c=1\)的單位制下，洛侖茲變換為
\begin{align}
x'&=\frac{x-vt}{\sqrt{1-v^2}},&
t'&=\frac{t-vx}{\sqrt{1-v^2}}.
\label{eq:sr1-lorentz-c1}
\end{align}

\begin{figure}[htbp]
\centering
\includegraphics[width=0.90\textwidth]{lecture_01_figure_08.png}
\caption{在光速設為一的單位制下，完整的洛侖茲變換。}
\lectureframe{01:02:05}
\end{figure}

\subsection{還原一般單位}

利用量綱分析即可還原一般單位。分母只能依賴無量綱比值\(v^2/c^2\)；時間方程中的\(vx\)則必須改為\(vx/c^2\)：
\begin{align}
x'&=\gamma(x-vt),&
t'&=\gamma\left(t-\frac{vx}{c^2}\right),
\label{eq:sr1-lorentz-dimensional}\\
\gamma&=\frac{1}{\sqrt{1-v^2/c^2}}.
\label{eq:sr1-gamma}
\end{align}
逆變換為
\begin{align}
x&=\gamma(x'+vt'),&
t&=\gamma\left(t'+\frac{vx'}{c^2}\right).
\label{eq:sr1-lorentz-inverse-dimensional}
\end{align}

\begin{figure}[htbp]
\centering
\includegraphics[width=0.90\textwidth]{lecture_01_figure_09.png}
\caption{還原\(c\)後的洛侖茲變換公式：\(v^2/c^2\)決定比例因子，\(vx/c^2\)則使位置混入時間。}
\lectureframe{01:05:01}
\end{figure}

當\(v\ll c\)時，
\begin{equation}
\gamma=1+\frac{1}{2}\frac{v^2}{c^2}+O\!\left(\frac{v^4}{c^4}\right),
\end{equation}
變換便簡化為
\begin{equation}
x'\approx x-vt,
\qquad
t'\approx t.
\end{equation}
在地球上常見的速度下，這些比值通常都很小，這就解釋了牛頓運動學為何如此成功。接近\(c\)時，修正則會變得很大。

這裡選擇沿\(x\)軸運動。若要符合互易性，橫向量尺在兩個慣性參考系中的測量結果就不能不同，因此
\begin{equation}
y'=y,
\qquad
z'=z.
\end{equation}

\section{測量運動中的桿長}

長度是指在進行測量的參考系中，同一時刻測得的兩端點間距。這不等於照片中物體看起來的視角大小，因為光傳播所需的時間會帶來另一個獨立的光學問題。

設一根桿靜止於不帶撇座標的參考系，兩端點位於
\begin{equation}
x=0,
\qquad
x=1.
\end{equation}
帶撇座標的觀測者在同一個帶撇同時面上測量兩端點的事件，例如取\(t'=0\)。在遠端\(x=1\)處，
\begin{equation}
t'=0
\quad\Longrightarrow\quad
t=\frac{vx}{c^2}.
\end{equation}
暫時回到\(c=1\)的單位制，這個事件的\(t=v\)為：
\begin{equation}
x'=\frac{x-vt}{\sqrt{1-v^2}}
   =\frac{1-v^2}{\sqrt{1-v^2}}
   =\sqrt{1-v^2}.
\end{equation}

\begin{figure}[htbp]
\centering
\includegraphics[width=0.90\textwidth]{lecture_01_figure_10.png}
\caption{桿靜止於不帶撇座標參考系時，其世界線為\(x=0\)和\(x=1\)；運動中的觀測者必須在同一個\(t'\)切面上比較兩者。}
\lectureframe{01:11:14}
\end{figure}

\begin{figure}[htbp]
\centering
\includegraphics[width=0.90\textwidth]{lecture_01_figure_11.png}
\caption{帶撇座標參考系在同一時刻測量兩端點，得到收縮因子\(\sqrt{1-v^2}\)（單位制為\(c=1\)）。}
\lectureframe{01:16:03}
\end{figure}

還原單位，並令靜止長度為一般值\(L_0\)，則
\begin{equation}
L=\frac{L_0}{\gamma}
 =L_0\sqrt{1-\frac{v^2}{c^2}}.
\label{eq:sr1-length-contraction}
\end{equation}

\begin{classroomqa}
\audiencequestion{如果桿的遠端位於\(x=1\)，它的座標不也應該是\(x'=1\)嗎？}\lecturerresponse{不是。桿靜止於不帶撇座標參考系，所以\(x=0\)和\(x=1\)標示的是它在該參考系中的固定端點。遠端的帶撇座標會隨時間改變；洛侖茲變換算出的，正是測量切面\(t'=0\)上該端點的座標。\lecturetimestamp{01:14:24}}\end{classroomqa}

\begin{classroomqa}\audiencequestion{為什麼每位慣性觀測者都能一致地說，另一位觀測者的桿比較短？}\lecturerresponse{他們測量的是不同的端點事件組合。不帶撇座標的觀測者要求\(t\)相等；帶撇座標的觀測者要求\(t'\)相等。由於同時性取決於參考系，兩人測得的並不是同一組端點事件。因此，互易性不會造成矛盾。\lecturetimestamp{01:16:28}}\end{classroomqa}

值得明確說明實際的測量方式。在測量參考系中沿途放置同步時鐘和助手。在指定的當地時間，每位助手記錄從身旁經過的端點位置。將兩個同時記錄的位置相減，就得到桿在該參考系中的長度。這項測量不取決於桿看起來是什麼樣子。

\begin{classroomqa}\audiencequestion{當\(c=1\)時，單位長度的桿一定要長一光秒嗎？}\lecturerresponse{只有時間以秒為單位時才是如此。也可以用公尺測量長度，並以光通過一公尺所需的時間來測量時間。設定\(c=1\)只是把時間與長度的單位連結起來，並未規定必須先選擇哪一種物理單位。\lecturetimestamp{01:21:17}}\end{classroomqa}

\section{運動時鐘與雙生子問題}

現在考察一個靜止於帶撇座標參考系的時鐘。沿著它的世界線，
\begin{equation}
x'=0.
\end{equation}
比較它讀數為\(t'=1\)時的事件與原點事件。逆變換給出
\begin{equation}
t=\gamma\left(t'+\frac{vx'}{c^2}\right)=\gamma.
\end{equation}
因此，運動時鐘上每經過一個時間單位，不帶撇座標參考系中的座標時間就經過\(\gamma\)個單位：
\begin{equation}
\Delta t=\gamma\,\Delta\tau,
\qquad
\Delta\tau=\frac{\Delta t}{\gamma}.
\label{eq:sr1-time-dilation}
\end{equation}

\begin{figure}[htbp]
\centering
\includegraphics[width=0.90\textwidth]{lecture_01_figure_12.png}
\caption{沿\(x'=0\)攜帶的時鐘，在同一對事件之間記錄的流逝時間，少於靜止座標時間：\(t=\gamma t'\)。}
\lectureframe{01:26:42}
\end{figure}

這種效應適用於所有物理時鐘，而不只是標示為時鐘的機械裝置。腕錶、原子躍遷、新陳代謝和生理老化，都遵循相同的固有時。在紐約街頭買來的廉價手錶，也不會因為比勞力士差就能逃脫相對論；它或許走時不準，但無論採用什麼機制，都仍受相同運動學支配。

\begin{classroomqa}
\audiencequestion{如果每位慣性觀測者都看到對方的時鐘變慢，為什麼旅行的雙生子回來時會比較年輕？}\lecturerresponse{往返旅行並非比較兩條在整段旅程中都等價的慣性世界線。旅行者必須折返，從一個慣性參考系轉到另一個；留在家中的雙生子則始終沿著同一條慣性世界線運動。折返打破了兩者的對稱性，因此累積的固有時可以不同。\lecturetimestamp{01:26:37}}\end{classroomqa}

\begin{classroomqa}\audiencequestion{年齡差異是由折返時的加速度本身造成的嗎？}\lecturerresponse{加速度標示出對稱性被打破的位置，但年齡差異是由沿完整路徑積分固有時而得。折返所需的時間可以很短，也可以把旅程視為在不同慣性運動階段之間切換；即使如此，兩條不同的時空路徑仍會累積不同的固有時。\lecturetimestamp{01:29:20}}\end{classroomqa}

\section{不變的幾何}

若用歐幾里得幾何來解讀，時間膨脹圖看起來就很違反直覺。傾斜的路徑看起來像斜邊，累積的固有時卻比垂直路徑短。問題的解答不在代數把戲，而在於時空具有不同的度規。

在歐幾里得平面中，旋轉會改變\(x\)和\(y\)，但會保留
\begin{equation}
x'^2+y'^2=x^2+y^2.
\label{eq:sr1-euclidean-invariant}
\end{equation}
因此很自然會檢驗洛侖茲變換是否保留\(t^2+x^2\)。答案是否定的。在\(c=1\)的單位制下，
\begin{align}
t'^2
&=\frac{t^2-2vtx+v^2x^2}{1-v^2},\\
x'^2
&=\frac{x^2-2vtx+v^2t^2}{1-v^2}.
\end{align}
將兩式相加會得到不對的係數；改為相減，就能消去混合項，得到
\begin{align}
t'^2-x'^2
&=
\frac{t^2-2vtx+v^2x^2-x^2+2vtx-v^2t^2}{1-v^2}\\
&=t^2-x^2.
\label{eq:sr1-interval-invariance}
\end{align}

\begin{figure}[htbp]
\centering
\includegraphics[width=0.90\textwidth]{lecture_01_figure_13.png}
\caption{將變換後的座標平方相減，混合項便會消去，並顯示出不變量\(t^2-x^2\)。}
\lectureframe{01:38:07}
\end{figure}

座標是隨參考系改變的分量，但以下組合
\begin{equation}
\Delta s^2=c^2\Delta t^2-\Delta x^2
\label{eq:sr1-spacetime-interval}
\end{equation}
在每個慣性參考系中都相同。若空間維度更多，
\begin{equation}
\Delta s^2
=c^2\Delta t^2-\Delta x^2-\Delta y^2-\Delta z^2.
\end{equation}

\begin{classroomqa}
\audiencequestion{為什麼可以不採用斜邊比較長這條一般幾何規則？}\lecturerresponse{那條規則適用於正定的歐幾里得度規。這裡的座標是透過尺、鐘和光訊號等操作方式定義的，而座標變換保留的是平方差。閔可夫斯基幾何不是為了避開悖論而預先假設的；它是從這些測量中推導出來的。\lecturetimestamp{01:43:21}}\end{classroomqa}

\subsection{固有時}

考慮時鐘世界線上的兩個事件，並選取時鐘的靜止參考系。在該參考系中，\(\Delta x'=0\)，所以不變性給出
\begin{equation}
c^2\Delta\tau^2
=c^2\Delta t^2-\Delta x^2,
\end{equation}
也就是
\begin{equation}
\Delta\tau
=\sqrt{\Delta t^2-\frac{\Delta x^2}{c^2}}.
\label{eq:sr1-proper-time}
\end{equation}
這就是時鐘在這兩個事件之間實際記錄到的時間。

\begin{figure}[htbp]
\centering
\includegraphics[width=0.90\textwidth]{lecture_01_figure_14.png}
\caption{在時鐘的靜止參考系中，空間間隔為零，因此不變間隔就是固有時的平方。}
\lectureframe{01:41:28}
\end{figure}

\begin{classroomqa}\audiencequestion{固有時應該是\(t^2-x^2/c^2\)，還是它的平方根？}\lecturerresponse{這個平方差等於固有時的平方。若間隔是從共同原點量起，時鐘實際讀取的流逝時間是\(\tau=\sqrt{t^2-x^2/c^2}\)。最初表述中漏掉的平方根必須補上。\lecturetimestamp{01:51:02}}\end{classroomqa}

\subsection{類時、類光與類空間隔}

兩個不同事件之間的間隔也可能為零。在\(c=1\)的單位制下，
\begin{equation}
\Delta t^2-\Delta x^2=0
\quad\Longleftrightarrow\quad
\Delta x=\pm\Delta t.
\end{equation}
這樣的事件由光線連接。

\begin{figure}[htbp]
\centering
\includegraphics[width=0.90\textwidth]{lecture_01_figure_15.png}
\caption{零間隔事件位於\(x=\pm t\)上；時空間隔與光樣路徑的固有時皆為零。}
\lectureframe{01:46:40}
\end{figure}

假設光能帶著一般意義上的時鐘，那麼從發射到吸收之間，時鐘不會累積任何固有時。這並不表示光有一個有效的靜止參考系，而是說每一段零間隔路徑的固有時都為零。

光線將時空分成三種不變的類別：
\begin{align}
c^2\Delta t^2-\Delta x^2&>0
&&\text{timelike},\\
c^2\Delta t^2-\Delta x^2&=0
&&\text{null},\\
c^2\Delta t^2-\Delta x^2&<0
&&\text{spacelike}.
\end{align}
物理時鐘可以連接由類時分隔的事件，而正的不變間隔決定時鐘的固有時。亞光速物體無法連接由類空分隔的事件。對於一對類空事件，可以選擇一個參考系，使兩個事件同時發生，並定義固有距離
\begin{equation}
\ell_{\rm proper}
=\sqrt{\Delta x^2-c^2\Delta t^2}.
\end{equation}

\begin{figure}[htbp]
\centering
\includegraphics[width=0.90\textwidth]{lecture_01_figure_16.png}
\caption{光錐將由類時分隔、可由時鐘連接的事件對，與由類空分隔、無法由任何亞光速物體連接的事件對區分開來。}
\lectureframe{01:48:47}
\end{figure}

\begin{classroomqa}\audiencequestion{什麼情況下應稱這個不變量為固有時，什麼情況下又應稱為固有長度？}\lecturerresponse{對於一對類時事件，存在一個參考系，使兩個事件發生在同一地點；該間隔給出固有時。對於一對類空事件，存在一個參考系，使兩個事件發生在同一時刻；符號相反的不變量給出兩者間的固有距離。類光間隔則是兩者之間的界線。\lecturetimestamp{01:47:22}}\end{classroomqa}

\begin{classroomqa}\audiencequestion{當單位沒有選成\(c=1\)時，\(c\)應該如何表示？}\lecturerresponse{不變量是\(c^2\Delta t^2-\Delta x^2\)，量綱為長度平方。將它除以\(c^2\)，得到\(\Delta\tau^2=\Delta t^2-\Delta x^2/c^2\)，量綱為時間平方。\lecturetimestamp{01:49:56}}\end{classroomqa}

\section{洛侖茲對稱性及其歷史意義}

\begin{classroomqa}\audiencequestion{洛侖茲變換可以視為旋轉嗎？}\lecturerresponse{可以，但不是普通的歐幾里得旋轉，而是保留\(t^2-x^2\)的雙曲旋轉。從歷史角度來說，可以把它描述成轉過一個虛角；現代則以快度作為實數參數。\lecturetimestamp{01:51:40}}\end{classroomqa}

愛因斯坦的思路與電磁學密切相關。根據他後來的回憶，他還是少年時，就想過如果自己追著光波跑會發生什麼事。按照伽利略速度相減，光波看起來會靜止下來；但麥克斯韋方程組並沒有描述真空中有一種靜止振盪的光波。這個衝突指向的是運動學，而不是尺具因運動而產生的某種特殊力學效應。

\begin{classroomqa}\audiencequestion{愛因斯坦的論證主要是回應邁克爾遜--莫雷等實驗嗎？}\lecturerresponse{這些實驗在歷史脈絡中確實重要，但愛因斯坦本人的說法著重於麥克斯韋方程組的結構，以及追趕光束的思想實驗。至於他對特定干涉儀實驗結果了解多少、是否依賴這些結果，歷史學界仍有爭論。因此，清楚的物理要點是麥克斯韋方程組與伽利略變換之間的對稱性衝突。\lecturetimestamp{01:52:04}}\end{classroomqa}

麥克斯韋方程組在洛侖茲變換下保持不變，在伽利略變換下則不然。洛侖茲和菲茨傑羅德早已找到相關的收縮公式，但他們仍保留以太，並將收縮解釋為物體穿過以太時受運動影響而產生的動力學形變。依照這種看法，該效應可能只是近似的物質反應，還會有其他修正。

愛因斯坦的關鍵一步，是將相對性原理與光速不變確立為精確的運動學原理。於是，長度收縮和時間膨脹適用於所有合適的尺與鐘，因為它們反映的是慣性參考系之間的幾何關係，而不是以太施加的特殊作用力。接下來要做的，是在這套幾何中描述粒子的運動。
